\clearpage
\newpage
\addcontentsline{toc}{chapter}{ABSTRACT}

\begin{center}
  \MakeUppercase{\@title}\\
  \bigskip
  \MakeUppercase{\@author}\\
  \bigskip
  \textbf{ABSTRACT}\\
\end{center}

\bigskip
Cyclostationary signal processing (CSP) provides the ability to estimate received waveforms' statistical features blindly.  Quadrature amplitude modulated (QAM) waveforms, when filtered by the square-root-raised cosine (SRRC) pulse shape function, have cyclic features that CSP can exploit to detect waveform parameters such as symbol rate (SR) and center frequency (CF).  The estimation of these SR-CF pairs enables a cognitive radio (CR) to perform spectrum sensing techniques such as spectrum sharing and interference mitigation.  Here, we investigate a field-programmable gate array (FPGA) application of a blind symbol rate-center frequency estimator.  First, this study provides a background on the theory behind the cyclic spectral density function (CSD), spectral correlation analyzers (SCA), and spectrum sensing.  Following this is a discussion on the motivation for CubeSat spectrum sensing.  An SCA implementation for low-memory devices, such as FPGA-based CubeSat, is then describes.  The paper concludes by reporting the performance characteristics of the newly developed streaming-based SCA.
\vfill

